% Propietario conservado: /Users/ruben/Documents/New project/output/REVISION_ARGUMENTAL_APERTURAS_CIERRES_20260915/VI/delivery/MANUSCRIPT_VI_ES_EN_COMPLETE/source_es/sections/14_generacion_regional.tex
\section{Initial conditions and regional generation}
\label{vi:vi:reg:seccion}\label{vi:sec:extension}\label{vi:sec:generacion}

An APP position does not by itself contain the initial condition of its
dynamics. The latter includes the direction of transport and distinguishes the cursors
of the two sheets. Before classifying particles, this domain must be
constructed: it produces the regional records that
will enter the electronic coordinate, action, and mass
readers. The operations in this section receive neither a physical species nor
a metrological value as an initial condition.

\subsection{Two oriented cursors on the complete support}

We use indices \(I_9=\{0,\ldots,8\}\), whose positive label is
\(i+1\), and the cardinal directions \(\mathcal D=\{N,E,S,O\}\).
The domain of one cursor and the joint domain are
\begin{equation}
 \mathcal S=I_9^2\times\mathcal D,\qquad
 \Omega_{\mathrm{seed}}=\mathcal S_+\times\mathcal S_\times,
 \qquad |\mathcal S|=324,\quad
 |\Omega_{\mathrm{seed}}|=104\,976.
 \label{vi:eq:semillas}
\end{equation}
The subscript \(\mathrm{seed}\) retains the domain designation used in
the antecedents; mathematically, these are oriented initial
conditions. The Cartesian product includes all ordered pairs.
The additive sheet has one cursor and the multiplicative sheet another; a
subsequent coincidence between their readings is a result of the traversal.

During the \(54\) instants of the initial emission, the phase selector
repeats \((+1,-1,0)\). In the positive regime,
\(\rho_9(i+j+2)\) is first evaluated at the additive cursor, and that
cursor is then transported by one position. In the negative regime,
\(\rho_9((i+1)(j+1))\) is evaluated at the multiplicative cursor, after which
the second cursor is transported. The null regime records a neutral event. The
directions are reversed after instants \(27\) and \(54\).
The latter reversal belongs to the terminal state although it does not alter
readings that have already been taken.

The multiplicative observable uses the exponent chart of the group of
units of \(\ZZ/9\ZZ\):
\[
 \ell_9(1)=0,\quad\ell_9(2)=1,\quad\ell_9(4)=2,\quad
 \ell_9(8)=3,\quad\ell_9(7)=4,\quad\ell_9(5)=5.
\]
In the sector \(\{3,6,9\}\), this observable takes the value zero.
Membership in the radical sector remains present in the state; assigning
that value to a scalar reading does not identify the radical with the multiplicative
identity.

\subsection{Finite formula for the window signatures}

The six windows each contain nine instants. Each sheet is evaluated
three times per window. A closed formula reconstructs the
signatures without consulting a prior list. For a cursor with origin
\(p=(i,j)\) and direction \(d\), let
\[
 f(k)=\begin{cases}k,&0\le k<9,\\18-k,&9\le k<18,
 \end{cases}\qquad p_k=p+f(k)\nu_d\pmod9.
\]
The index \(k\) counts only the evaluations of that sheet. The
reversal after nine of its evaluations explains the second
expression for \(f\). Writing \(p_k=(i_k,j_k)\), for
\(m=0,\ldots,5\) we obtain
\begin{align}
 S_m(p,d)&=\sum_{k=3m}^{3m+2}\rho_9(i_k+j_k+2),\notag\\
 E_m(p,d)&=\sum_{k=3m}^{3m+2}
       \ell_9\!\left(\rho_9((i_k+1)(j_k+1))\right),\notag\\
 s_m&=[S_m]_{10},\qquad e_m=[E_m]_{10}.
 \label{vi:vi:reg:firmas}
\end{align}
The integer evaluation and its division quotient remain stored in
the record. The preceding signatures are the decimal projections
specified by the emitter.

Since each window has \(Z_m=3\) neutral events, the rule
coupling the two signatures is
\begin{align}
 d_{m,1}&=[S_m]_{10},\qquad
 d_{m,2}=[S_m+E_m]_{10},\notag\\
 d_{m,3}&=[3S_m+5E_m+7Z_m]_{10},\notag\\
 U_m&=100s_m+10[s_m+e_m]_{10}+[3s_m+5e_m+1]_{10}.
 \label{vi:eq:emisor-seis}
\end{align}
The final term \(1\) is the reduction of \(7Z_m=21\) modulo
ten. The coefficients and calendar are stated in the emission
law; physical magnitudes do not appear among its arguments.
The oriented ternary projection is
\begin{equation}
 \Psi(U_0,\ldots,U_5)=([-U_0]_3,\ldots,[-U_5]_3)
 \in\FF_3^6.
 \label{vi:eq:psi-catalogo}
\end{equation}
The balanced identification of TRIT is then applied. The minus
sign belongs to this chart and must be transported when it changes.

\begin{proposition}[Factorisation and census of the emission]
\label{vi:vi:reg:censo}
The initial conditions produce \(18\) additive signatures and \(26\)
multiplicative signatures. Their coupling is injective and has \(468\)
emissions. The image of \(\Psi\) contains \(243\) words in the
ambient space of \(729\) elements.
\end{proposition}
\begin{proof}
The hundreds digits of \(U_m\) recover \(s_m\). If \(b_m\) is the
tens digit, \(e_m=[b_m-s_m]_{10}\); the units digit
checks the third congruence. Hence no distinct pair
of signatures has the same emission.

To obtain the census, \eqref{vi:vi:reg:firmas} is evaluated at
\(i,j=0,\ldots,8\) and the four directions. Appendix
\ref{vi:vi:reg:tablas} contains all the resulting signatures. Each
additive signature has \(18\) preimages. Among the multiplicative signatures,
\(24\) have \(8\), one has \(24\), and one has \(108\).
These sums are respectively \(18\cdot18=324\) and
\(24\cdot8+24+108=324\). Since the two cursors are
independent, the multiplicities multiply. This gives
\(432\) emissions of weight \(144\), \(18\) of weight \(432\), and
\(18\) of weight \(1944\), with sum
\[
 432\cdot144+18\cdot432+18\cdot1944=104\,976.
\]
The six reductions of \eqref{vi:eq:psi-catalogo} are applied to each pair
of the \(18\) and \(26\) signatures published in the appendix. The
complete orbit decomposition included there gives
\(38\cdot6+5\cdot3=243\) different images. All the
operations of this enumeration are the finite sums, products and residues
of the preceding formulas.
\end{proof}

The finite proof does not identify exhaustiveness of the domain with
unlimited depth. The initial conditions have been enumerated on
a fixed support. Depth is constructed through compatible
prolongations preserving those conditions and their transports.

\subsection{Selection of regions through oriented fibres}

For \(w=(w_1,\ldots,w_6)\), define
\[
 r(w)=(w_3,w_1,w_2,w_5,w_6,w_4),\qquad
 s(w)=(w_4,w_5,w_6,w_1,w_2,w_3).
\]
The identities \(r^3=s^2=I\), \(srs=r^{-1}\) realise
\(D_3\). Permuting the six components of \(U\) in the same way
commutes with \(\Psi\). The preceding enumeration verifies that the
orbits preserve the multiplicities of their fibres.

Let \(f(w)\) be the number of emissions over \(w\), \(m(w)\)
its total multiplicity, and \(\nu(w)=(n_0,n_1,n_2)\) its ternary
census. The additive signature preserves the diagonal class
\(d=[i+j]_9\) of the origin and the sense \(ES\) or \(NO\).
We write \(\mathcal D(\mathcal O)\) for the set of these
classes over an orbit.

\begin{proposition}[Distinguished regions of the catalogue]
\label{vi:vi:reg:seleccion}
There is a unique orbit of six words with five emissions per
word and multiplicities \(432+4\cdot144\). Among the orbits
of six words satisfying
\[
 f(w)=1,\quad m(w)=144\quad(w\in\mathcal O),\qquad
 \mathcal D(\mathcal O)=\{0,3,6\},
\]
the census \((2,3,1)\) and the census \((2,2,2)\) each select
a unique orbit. The mark \(d=6\) with orientation
\(ES\) determines the representatives
\[
 w^{\rm cl}=010211,\qquad w^{\rm pr}=201101,
 \qquad w^{\rm au}=121200.
\]
\end{proposition}
\begin{proof}
The complete table of the \(43\) orbits in
\ref{vi:vi:reg:tablas} allows the predicates to be applied. The closure
orbit is the only one with a fibre of size five and weight \(1008\);
its five weights are those stated. The diagonal condition and the
simple fibre of weight \(144\) leave six orbits. Their censuses
select respectively one with \((2,3,1)\) and one with
\((2,2,2)\). Applying the six permutations of \(D_3\) and
checking the diagonal class and orientation in
\eqref{vi:vi:reg:firmas} fixes the three representatives displayed.
Uniqueness uses all the conditions: if the set of
diagonals is omitted, four balanced orbits with simple fibre
and weight \(144\) appear.
\end{proof}

The designations closure, propagation, and self-scaling apply
first to these regions. The next section introduces their
evaluation characters. In particular, the multiplicative signature
\(555555\) distinguishes a transverse closure emission, whereas
\((5,5)\) is a position of the electronic support. The two
constructions are related through subsequent operations, not
through identity of their digits.

\subsubsection{Mirror involution of the regional components}
\label{vi:mono-ampl:regiones}

On the regional block
\(B=(u^{\rm cl},u^{\rm pr},u^{\rm au})\in(\FF_3^6)^3\),
the interchange is
\[
 \mathfrak cB=(u^{\rm pr},u^{\rm cl},u^{\rm au}).
\]
Its square is the identity and it preserves the self-scaling component.
If \(q=u^{\rm cl}+u^{\rm pr}+u^{\rm au}\),
\(a=u^{\rm cl}+u^{\rm pr}-u^{\rm au}\), and
\(d=u^{\rm cl}-u^{\rm pr}\), the block is recovered by
\[
 u^{\rm au}=2(q-a),\quad s=q-u^{\rm au},\quad
 u^{\rm cl}=2(s+d),\quad u^{\rm pr}=2(s-d).
\]
The proof uses \(2^{-1}=2\) in \(\FF_3\). The interchange
preserves \(q,a,s\) and reverses the sign of \(d\). Thus
the symmetric aggregate and the oriented difference jointly preserve
the two components that the aggregate alone does not distinguish.
This regional involution and the reversal of a cardinal route
act on different domains. The fixed component may
evolve during transport even though it is fixed by the interchange.

\subsection{Calendar, monodromy, and memory}

The finite calendar consists of twelve nonadic windows. Its
group sequence is
\begin{equation}
 0\longrightarrow C_{12}\xrightarrow{[b]\mapsto[9b]}C_{108}
 \xrightarrow{[t]\mapsto[t]_9}C_9\longrightarrow0.
 \label{vi:eq:extension108}
\end{equation}
For \(t=9b+r\), \(0\le r<9\), composition contains
the transport quotient
\[
 (b,r)+(b',r')=
 \left(b+b'+\left\lfloor\frac{r+r'}9\right\rfloor,
       [r+r']_9\right),
\]
with the first coordinate reduced modulo twelve. The projection is
exact because its kernel consists of the multiples of nine. It admits no
homomorphic section: such a section would give \(C_{108}\cong C_{12}\times C_9\),
but the respective exponents are \(108\) and \(36\).

Monodromy describes the action of one circuit of the phase base
on its fibre. In the dynamics with memory, the nonadic holonomy
\(\Gamma_9\) preserves phase and advances depth. Its
projection onto phase and counter satisfies
\[
 (r,n)\longmapsto(r,n+1).
\]
This formula is a reduced representation of the action, not a
definition eliminating the rest of the state. On bidirectional
histories the counter can be extended to \(\ZZ\); on
histories starting at depth zero, backward motion is
defined only where the predecessor is preserved. The calendar can
return after \(108\) transitions while preserving a memory
different from the initial one. This distinction will be essential when
the particle is constructed as persistence rather than repetition
of a visible coordinate.

\subsection{From records to prolongation operators}

The linear publication operators are determined on the
transport transitions. If \(x_j^\chi\) is a regional state
and \(\Pi_6\) its six-component publication, the dependence is
\[
 b_j^\chi=\Pi_6(x_j^\chi),\qquad
 x_{j+1}^\chi=U_{9j+9}\cdots U_{9j+1}(x_j^\chi).
\]
The origin matrices \(X_a\) and destination matrices \(Y_a\), in
row vectors of \(\FF_3^6\), are
\[
\begin{array}{c|c|c|c}
X_0&Y_0&X_1&Y_1\\\hline
010211&012222&010211&002111\\
012222&010211&002111&110221\\
201101&121221&102011&012222\\
121221&102011&012222&102011\\
121200&112202&121020&010210\\
112202&121020&010210&010200
\end{array}
\]
We have \(\det X_0=1\), \(\det X_1=-1\) in \(\FF_3\).
The unique solution of \(X_aL_a=Y_a\) is \(L_a=X_a^{-1}Y_a\).
Concatenation of blocks published in the calendar
\((L_0,L_0,L_1,L_1)\) produces lengths
\(6,12,18,24,30\). The boundary of length \(36\) preserves
the terminal and the prolongation relation. These matrices
represent transitions: linear inversion proves their uniqueness
for the given record, while generation of the states
belongs to the effective composition \(U_t\) of the TPK.

\subsection{Scope of the regional construction in the study of masses}

The census has separated initial conditions, emissions, regions, and
operators representing their transitions. The next stage
evaluates the regional coordinates and composes the dodecaphase record.
Later, the particle atlas will be constructed on the same
support but will use another classification. This continuity preserves
a common provenance without turning a numerical fibre into a
species through mere coincidence of cardinalities or labels.
